\begin{vidu}%[1P5-1.2K][Loai1]
	Lượng chất (số mol) chứa trong 924 g khí $CO_2$ là
	\choice
	{4,2 mol}
	{2,1 mol}
	{\True 21 mol}
	{42 mol}
	\loigiai{
		\begin{itemize}
			\item Phân tử gam của $CO_2$ là 44 gam
			\item Lượng chất chứa trong 924 g khí $CO_2$:
			\begin{align*}
				n=\dfrac{m_{CO_2}}{M_{CO_2}}=\dfrac{924}{44}=21\ mol. 
			\end{align*}
		\end{itemize}
	}
\end{vidu}
\begin{vidu}%[1P5-1.2K][Loai1]
	Một bình kín chứa $N=3,01.10^{23}$ phân tử khí Heli. Cho số Avogadro $N_A = 6,02.10^{23}\ mol^{-1}$. Khối lượng khí Helium chứa trong bình là
	\choice
	{\True 2 g}
	{4 g}
	{6 g}
	{8 g}
	\loigiai{
		Khối lượng khí Helium trong bình:
		\begin{align*}
			m_{He}=n.M_{He}=\dfrac{N}{N_A}M_{He}=\dfrac{{3,01.10}^{23}}{{6,02.10}^{23}}.4=2\ g.
		\end{align*}
	}
\end{vidu}

\begin{vidu}%[1P5-1.2K][Loai1]
	Trong một bình kín có chứa $10\ g$ khí oxygen. Cho số Avogadro $N_A = 6,02.10^{23}\ mol^{-1}$. Số phân tử khí oxygen chứa trong bình bằng
	\choice
	{\True $19.10^{22}$ phân tử}
	{$1,9.10^{22}$ phân tử}
	{$19.10^{23}$ phân tử}
	{$1,9.10^{21}$ phân tử}
	\loigiai{
		Số phân tử khí chứa trong bình: $N = \dfrac{m}{M}.N_A=\dfrac{10}{32}.6,02.10^{23}= 19.10^{22} $ phân tử.
	}
\end{vidu}

\begin{vidu}%[1P5-1.2K][Loai1]
	Một bình kín chứa $N= 3,01.10^{23}$ phân tử khí heli. Cho số Avogadro $N_A = 6,02.10^{23}\ mol^{-1}$. Biết nhiệt độ khí là $0\doC$ và áp suất khí trong bình là 1 atm ($1,013.10^5\;Pa$). Thể tích của bình là
	\choice
	{22,4 lít}
	{\True 11,2 lít}
	{5,6 lít}
	{2,24 lít}
	\loigiai{
		\begin{itemize}
			\item Trong điều kiện nhiệt độ và áp suất như trên (ĐKTC), thể tích của 1 mol He là $V_0=22,4$ lít. 
			\item Vì lượng khí He trong bình chỉ là $n=\dfrac{N}{N_A}=0,5$ mol nên thể tích của bình là:
			$V=n.V_0=11,2$ lít.
		\end{itemize}
	}
\end{vidu}
\begin{vidu}%[1P5-1.2K][Loai1]
	Trong một bình có thể tích $2$ lít chứa $10\ g$ khí oxygen. Cho số Avogadro $N_A = 6,02.10^{23}\ mol^{-1}$. Khối lượng riêng của khí là
	\choice
	{0,5 $kg/m^3$}
	{\True 5 $kg/m^3$}
	{0,05 $kg/m^3$}
	{0,005 $kg/m^3$}
	\loigiai{
		Khối lượng riêng của khí: $D = \dfrac{m}{V}=\dfrac{10.10^{-3}}{2.10^{-3}}=5\ kg/m^3 $.
	}
\end{vidu}
\begin{vidu} %[1P5-1.2Y][Loai1]
	Biết khối lượng của một mol nước là 18 g và 1 mol có $N_A= 6,02. 10^{23}$ phân tử. Số phân tử trong 2 g nước là
	\choice
	{$4. 10^{21}$ phân tử}
	{$1,8. 10^{20}$ phân tử}
	{$3,24. 10^{24}$ phân tử}
	{\True $6,69. 10^{22}$ phân tử}
	\loigiai{
		Số phân tử trong 2 g nước là: $N=n.N_A=\dfrac{m}{M}.N_A=\dfrac{2}{18}.6,02.10^{23}=6,69.10^{22}$ phân tử.
	}
\end{vidu}

\begin{vidu}%[1P5-1.2K][Loai1]
	Một bình kín chứa $5,6.10^{23}$ phân tử khí nitrogen. Cho số Avogadro $N_A = 6,02.10^{23}\ mol^{-1}$. Biết khối lượng mol nitrogen là 28,0134 g/mol. Khối lượng khí nitrogen trong bình là
	\choice
	{13 g}
	{\True 26 g}
	{1,3 g}
	{2,6 g}
	\loigiai{ 
		Ta có: $N=n.N_A=\dfrac{m}{M}.N_A\Rightarrow m=\dfrac{NM}{N_A}\approx 26\ g$.
	}
\end{vidu}
\begin{vidu} %[1P5-1.2B][Loai1]
	Biết khối lượng của 1 mol nước là $M = 18. 10^{- 3}$ kg và 1 mol có $N_A = 6, 02.10^{23}$ phân tử. Biết khối lượng riêng của nước là $D=10^3\ kg/m^3$. Số phân tử có trong $300\ m^3$ nước là
	\choice
	{$10, 03.10^{23}$ phân tử}
	{\True $10, 03.10^{24}$ phân tử}
	{$6, 7.10^{23}$ phân tử}
	{$6,7. 10^{24}$ phân tử}
	
	\loigiai{
		\begin{itemize}
			\item Số mol nước có trong $300\ cm^3$  nước là: 
			$n=\dfrac{m}{M}=\dfrac{D V}{M}=\dfrac{10^3.300.10^{ - 6}}{18.10^{ - 3}}=\dfrac{50}{3}\ mol.$
			\item Số phân tử có trong $300\ cm^3$ là $N=n.N_A= \dfrac{50}{3}.6, 02.10^{23}= 10, 03.10^{24}$ phân tử.
		\end{itemize}
		
	}
\end{vidu}